\section{Manuscript Naming Map}\label{apx:naming-map}
The manuscript uses family names for readability. Internal engineering stage names are mapped in Table~\ref{tab:naming-map} only for reproducibility.

\begin{table}[!htb]
\caption{Family naming used in the manuscript}
\label{tab:naming-map}
\centering
\small
\begin{tabular}{l l l}
\toprule
\textbf{Family name} & \textbf{Internal stage alias} & \textbf{Core idea} \\
\midrule
Uniform Baseline (UB) & Stage0 & uniform quantization with reconstruction-only objective \\
Adaptive ROI Student (ARS) & Stage1 & adaptive quantization with ROI supervision \\
Adaptive Distilled Student (ADS) & Stage2 & ARS + teacher-guided distillation \\
Multi-Scale Distilled Student (MS-ADS) & Stage3 & ADS + stronger multi-scale ROI branch \\
\bottomrule
\end{tabular}
\end{table}

\section{Concept-to-Implementation Traceability}\label{apx:traceability}
The main body presents theory and evidence. Table~\ref{tab:traceability} provides a conceptual traceability map without repository-specific details.

\begin{table}[!htb]
\caption{Traceability from objective terms to implementation components}
\label{tab:traceability}
\centering
\small
\begin{tabular}{p{0.23\textwidth} p{0.23\textwidth} p{0.43\textwidth}}
\toprule
\textbf{Objective term} & \textbf{Equation} & \textbf{Implemented component} \\
\midrule
Total objective composition & Eq.~\ref{eq:total-objective} & training loop objective aggregator \\
Reconstruction MSE & Eq.~\ref{eq:recon-loss} & reconstruction criterion module \\
Adaptive level mapping & Eq.~\ref{eq:level-map}--\ref{eq:rounded-level} & adaptive quantization module \\
Rate proxy variants & Eq.~\ref{eq:rate-global}--\ref{eq:rate-norm-bg} & rate regularization module with mode selection \\
Weighted ROI supervision & Eq.~\ref{eq:imp-bce}--\ref{eq:imp-posweight} & importance supervision module \\
ROI--BG margin regularization & Eq.~\ref{eq:sep-loss} & importance separation regularizer \\
Feature/logit distillation & Eq.~\ref{eq:distill-total} & teacher--student distillation module \\
End-to-end family pipeline & Eq.~\ref{eq:pipeline} & encoder, ROI branch, quantizer, decoder stack \\
\bottomrule
\end{tabular}
\end{table}

\section{Experimental Protocol Summary}\label{apx:protocol}
\begin{table}[!htb]
\caption{Protocol summary by model family}
\label{tab:protocol-summary}
\centering
\small
\begin{tabular}{l l l}
\toprule
\textbf{Family} & \textbf{Quantization strategy} & \textbf{Supervisory signals} \\
\midrule
UB & uniform latent quantization & reconstruction only \\
ARS & adaptive ROI-aware quantization & reconstruction + ROI supervision + rate regularization \\
ADS & adaptive ROI-aware quantization & ARS objective + distillation + ROI/BG separation \\
MS-ADS & adaptive ROI-aware quantization & ADS objective with multi-scale ROI branch variants \\
\bottomrule
\end{tabular}
\end{table}

Teacher-guided settings use a pretrained PointPillars-based detector teacher. Full-sweep experiments are 150 epochs, and quick diagnostic analyses are performed on a 16-frame subset.

\section{Data Artifacts and Figure Plan}\label{apx:artifacts}
The quantitative results in this draft are built from three artifact groups: (1) the global experiment ledger, (2) 16-frame quality diagnostics, and (3) native-vs-oracle decomposition summaries.

Figures are intentionally omitted in this revision. The next revision should include (i) precision/recall/IoU threshold sweeps, (ii) matched-budget bitrate--distortion plots, and (iii) qualitative panels showing input channels, ground-truth ROI, predicted importance, and thresholded ROI.
